\chapter{NEP}

\section{Preliminary plans}
\subsection{Second week}
\subsubsection{Goals for next week}
We are currently working to establish what thruster type we will use. Likely our next steps will be propellant and creating design targets. I will assign
\begin{itemize}
	\item Find NASA resources to create target
	\item Xeon(Hall, ion, RF)
	\item Krypton(hall or ion)
	\item Iodine(Hall, Ion)
	\item Argon(MHD or RF)
	\item Hydrogent(MHD, RF)
	\item Research Rover
	\item Research Atomic in Orbit
	\item Research Atomic at Lagrangian
\end{itemize}

\subsubsection{Goal}
Fly by, Rover, Atomic in Orbit, or Lagrandian

This document compares four candidate mission architectures for an
undergraduate nuclear-electric-propulsion (NEP) project centered on
the design and computational evaluation of a Hall-effect thruster:
(1) an atomic-clock spacecraft in Mars orbit, (2) an atomic-clock
spacecraft near the Mars--Sun first Lagrange point, (3) a spacecraft
performing a Mars flyby, and (4) a Mars rover mission.

The comparison separates mission utility from propulsion-system
suitability. These objectives overlap, but they are not equivalent.
A mission may have substantial scientific value while imposing
requirements that are poorly matched to the team's intended
propulsion study. Conversely, a comparatively modest mission may
provide a better engineering case if its trajectory, payload,
power demand, and operational lifetime can be translated into
well-defined thruster requirements.

The central conclusion is conditional. A Mars-orbit atomic-clock
mission has a strong connection to established concepts in deep-space
radiometric navigation and Mars communications infrastructure.
A clock near Mars--Sun L1 could provide different measurement
geometry, but its advantage must be demonstrated using a specified
navigation architecture rather than inferred from the Lagrange-point
location alone. A flyby is attractive for a focused transfer and
encounter study, whereas a rover mission offers an obvious surface
exploration objective but introduces atmospheric entry, descent,
landing, and surface-operation requirements that can obscure the
propulsion question.

No architecture can be declared the best Hall-thruster application
without a quantitative mission model. The most defensible selection
is the one for which the team can establish a traceable chain from
mission objectives to spacecraft mass, trajectory, electrical power,
propellant consumption, thrust, and thruster operating conditions.


% ============================================================
\paragraph{Scope, Assumptions, and Decision Criteria}
% ============================================================

\subparagraph{The actual decision}

The decision is not simply which mission is ``most useful.'' Four
distinct questions must be answered:

\begin{enumerate}
	\item \textbf{Mission utility:} What scientific, operational,
	      or infrastructure capability does the mission provide?

	\item \textbf{Trajectory compatibility:} Does the mission
	      benefit from continuous low thrust, or can it be accomplished
	      more simply using conventional chemical propulsion?

	\item \textbf{Thruster-design value:} Does the mission impose
	      sufficiently concrete requirements to guide channel geometry,
	      magnetic-field configuration, discharge power, propellant flow,
	      thermal management, and cathode integration?

	\item \textbf{Undergraduate feasibility:} Can the group evaluate
	      the mission with accessible trajectory tools, published
	      reference data, and a defensible simulation hierarchy without
	      needing to solve every spacecraft subsystem?
\end{enumerate}

These questions should be scored separately before a final mission
is selected. Otherwise, a mission may be chosen because it sounds
scientifically compelling even though its requirements do not
meaningfully constrain the thruster design.

\subparagraph{Assumed propulsion architecture}

The baseline assumption is a spacecraft using a Hall-effect
thruster powered by an electric power system ultimately associated
with a nuclear-electric-propulsion concept. The project focuses on
the thruster and its interaction with the mission, rather than on
designing a nuclear reactor.

The relevant spacecraft-level variables include

\begin{itemize}
	\item initial wet mass, \(m_0\);
	\item dry spacecraft mass, \(m_{\mathrm{dry}}\);
	\item propellant mass, \(m_p\);
	\item thruster electrical input power, \(P_{\mathrm{in}}\);
	\item thrust, \(T\);
	\item specific impulse, \(I_{\mathrm{sp}}\);
	\item total thruster operating time, \(t_{\mathrm{on}}\);
	\item required mission \(\Delta v\);
	\item mission duration and operational lifetime;
	\item payload power and spacecraft bus power;
	\item propellant type and storage requirements.
\end{itemize}

The analysis assumes that the team is evaluating a plausible mission
concept, not proposing a fully qualified flight system. A first-order
trajectory and propulsion analysis is appropriate, provided its
assumptions and limitations are explicit.

\subparagraph{Performance relations}

For an idealized electric thruster, the effective exhaust velocity is

\[
	v_e = g_0 I_{\mathrm{sp}},
\]

where \(g_0 \simeq 9.80665\ \mathrm{m\,s^{-2}}\).

The thrust is related to propellant mass flow by

\[
	T = \dot{m}v_e
	= \dot{m}g_0 I_{\mathrm{sp}},
\]

where \(\dot{m}\) is the propellant mass-flow rate. For a steady
operating point, an idealized jet-power relation is

\[
	P_{\mathrm{jet}}
	= \frac{1}{2}\dot{m}v_e^2
	= \frac{T v_e}{2}.
\]

If the total thruster efficiency is \(\eta_t\), then

\[
	P_{\mathrm{in}}
	\simeq \frac{P_{\mathrm{jet}}}{\eta_t},
\]

and therefore

\[
	T \simeq \frac{2\eta_t P_{\mathrm{in}}}{v_e}
	= \frac{2\eta_t P_{\mathrm{in}}}
	{g_0 I_{\mathrm{sp}}}.
\]

These equations are useful for preliminary sizing, but they do not
replace a real Hall-thruster performance map. Efficiency, discharge
voltage, mass utilization, divergence, cathode losses, magnetic
configuration, and thermal constraints all affect the actual
operating point.

The rocket equation gives a useful first-order estimate of
propellant requirements for a specified ideal velocity increment:

\[
	\Delta v
	= v_e \ln\left(\frac{m_0}{m_f}\right),
\]

where \(m_f\) is the final mass after propellant expenditure. This
equation is not, by itself, a trajectory solver. In particular, a
low-thrust trajectory must account for the direction and duration of
acceleration, changing spacecraft mass, planetary gravity, solar
gravity, and mission boundary conditions.

\subparagraph{What constitutes a good mission for this project?}

The most useful mission is one that constrains the thruster rather
than merely carrying it.

For example, a mission duration of several years may make high
specific impulse attractive because propellant mass is expensive.
However, increasing specific impulse at fixed electrical power
usually decreases thrust. A longer transfer may reduce the required
propellant or improve payload fraction while simultaneously increasing
the required spacecraft lifetime and total time spent operating the
thruster.

The mission therefore creates a trade between

\[
	\text{propellant mass},\quad
	\text{electrical power},\quad
	\text{thrust},\quad
	\text{transfer time},\quad
	\text{spacecraft lifetime}.
\]

A strong mission concept makes these competing requirements
quantifiable.

% ============================================================
\paragraph{Mission A: Atomic Clock in Mars Orbit}
% ============================================================

\subparagraph{Concept and potential purpose}

The spacecraft enters a bound orbit around Mars and carries a
high-stability atomic clock. Depending on the communications and
navigation architecture, the spacecraft could support one or more
of the following:

\begin{itemize}
	\item one-way radiometric tracking and spacecraft orbit
	      determination;
	\item autonomous or partially autonomous navigation;
	\item precision timing for a distributed Mars communications
	      or navigation system;
	\item radio-science measurements and improved orbit knowledge;
	\item timing and navigation experiments supporting future
	      Mars surface and orbital missions.
\end{itemize}

The clock is not, by itself, a navigation system. Its value depends
on the signals measured, the stability of the clock, the spacecraft's
orbit, the availability of external references, and the algorithms
used to infer position and velocity.

NASA's Deep Space Atomic Clock work explicitly examined one-way
radiometric tracking, autonomous navigation concepts, Mars-orbiter
orbit determination, and radio-science applications
\cite{DSACNavigation}. This provides a direct technical precedent
for an atomic-clock spacecraft in the Mars environment.

\subparagraph{Arguments in favor}

\textbf{A. The mission has a concrete operational context}

A clock-equipped Mars orbiter can be connected to a recognizable
operational need: accurately determining spacecraft position and
time while operating far from Earth.

Conventional deep-space tracking often relies on measurements
involving signals transmitted from Earth and received or retransmitted
by the spacecraft. Such tracking depends on ground assets, scheduling,
signal geometry, communications links, and the ability to interpret
the returned measurements.

A sufficiently stable onboard clock can support one-way
radiometric measurements. These can reduce dependence on a
round-trip measurement architecture and may permit more flexible
tracking and navigation operations. The potential benefit is not
that the clock removes the need for tracking, but that it changes
which measurements can be used and how often they can be collected.

The technical advantage is therefore conditional on the complete
measurement system. The clock must be sufficiently stable over the
relevant measurement interval, and the radio system must provide
measurements with suitable accuracy and geometry.

For a student project, this is valuable because the scientific
objective can be expressed in testable terms. For example:

\begin{quote}
	Determine whether a specified onboard clock stability and
	one-way radio-tracking architecture can achieve a target orbit
	determination accuracy for a spacecraft in a specified Mars orbit.
\end{quote}

That question can be studied using published clock stability
specifications, an orbital dynamics model, simulated radio
measurements, and an orbit-estimation algorithm.

\textbf{B. Mars already has an operational relay infrastructure}

Mars exploration is not limited to isolated landers. Orbiters have
served as communications relays between the Martian surface and
Earth. NASA describes the Mars Relay Network as a collection of
orbiters that transmit data from surface missions back to Earth
\cite{MarsRelayNetwork}.

This provides a plausible operational environment for studying
how precise timing could complement existing Mars infrastructure.
A clock-equipped spacecraft could potentially be designed as a
new node in a broader system rather than as an entirely isolated
scientific experiment.

There is an important distinction, however, between a communications
relay and a precision-navigation node. A spacecraft that forwards
data does not automatically provide the timing stability,
ranging capability, signal coverage, or orbit knowledge needed for
navigation. The mission architecture must specify which role the
clock-equipped spacecraft is expected to perform.

A defensible concept might therefore define the spacecraft as
a precision timing and navigation demonstrator that also supports
communications, rather than assuming that the clock alone creates
a complete Mars navigation service.

\textbf{C. The orbit provides repeated access to the Mars environment}

Unlike a flyby, a Mars orbiter can make repeated observations of
the same planetary environment. It can remain available for multiple
navigation experiments, repeated radio links, and comparisons
between predicted and observed spacecraft states.

This repeated access is useful for evaluating whether clock
performance improves orbit estimates over time. It also provides
a natural reason to consider long operational lifetimes, which may
be relevant to an electric-propulsion spacecraft.

The orbital design remains a major decision. A low Mars orbit,
high Mars orbit, eccentric orbit, and polar orbit differ in
coverage, orbital period, radiation environment, communications
geometry, and station-keeping requirements. The best orbit depends
on the mission's intended service, not simply on the fact that
the spacecraft is near Mars.

\textbf{D. The mission has a relatively clear connection to
	a Hall-thruster design}

An orbital mission allows the team to connect propulsion decisions
to a well-defined arrival condition. The spacecraft must depart
Earth, travel to Mars, and enter or achieve the desired bound orbit.

A low-thrust trajectory model can estimate the required thrust
history and propellant consumption. Those results constrain
the spacecraft's electrical power system and thruster operating
point.

The mission also creates a meaningful distinction between
thruster operation during transfer and operation during the
science phase. The thruster might be used primarily during
interplanetary cruise, during orbit insertion, for orbit changes,
or for some combination of these functions. The team should not
assume that the thruster must remain continuously active after
arrival.

A useful engineering study would compare at least two cases:

\begin{enumerate}
	\item chemical or other conventional injection into the
	      desired Mars orbit, followed by electric-propulsion
	      station-keeping or orbit adjustment;
	\item electric propulsion used for a substantial fraction
	      of the Earth-to-Mars transfer and, if required, arrival
	      orbit acquisition.
\end{enumerate}

This comparison determines whether the NEP system creates a
meaningful mass or mission-duration advantage rather than simply
adding complexity.

\textbf{E. The mission can be scoped around a measurable
	navigation objective}

A strong undergraduate project should define a baseline architecture
and then evaluate one change at a time.

For example, the team could compare a conventional tracking case
against a case with a specified onboard clock and one-way
radiometric tracking. Performance metrics could include

\begin{itemize}
	\item position and velocity estimation error;
	\item sensitivity to clock frequency drift;
	\item measurement availability;
	\item dependence on ground-station tracking;
	\item navigation performance under communication outages;
	\item sensitivity to uncertainty in the spacecraft's initial state.
\end{itemize}

The group would not need to build or validate a real atomic clock.
It would need to use defensible published clock specifications
and model their effect on navigation measurements.

\subparagraph{Arguments against}

\textbf{A. One clock-equipped orbiter is not a Mars GPS system}

The most important limitation is that a clock does not provide
position by itself. Navigation requires measurements that constrain
the spacecraft's state.

An onboard clock can support timing and one-way radiometric
tracking, but the achievable navigation accuracy depends on the
measurement architecture, signal geometry, orbit model, clock
noise, radio-system errors, and reference information.

A single spacecraft also does not automatically provide a global
positioning service to rovers or other orbiters. A Mars-local
navigation network would require a defined architecture with
appropriate signal sources, timing references, communications
links, coverage, and orbit knowledge. Depending on the design,
multiple spacecraft or additional ground and inter-spacecraft
measurements may be necessary.

Consequently, the claim that ``an atomic clock in Mars orbit will
provide Mars GPS'' is too strong unless a specific system is
designed and analyzed.

A better claim is:

\begin{quote}
	A clock-equipped Mars orbiter could demonstrate a component of
	a future Mars navigation architecture, provided that the
	measurement system and its navigation performance are specified.
\end{quote}

\textbf{B. The payload's benefit may not justify a dedicated
	spacecraft}

The clock could potentially be hosted on an existing orbiter or
included as a secondary payload. If so, a dedicated spacecraft
may not be necessary to demonstrate the clock's scientific
utility.

This creates a challenge for the mission justification. The team
must explain why a new NEP spacecraft is useful instead of placing
the clock on an existing or otherwise planned platform.

There are several possible answers: the mission could test a
new orbit, demonstrate autonomous navigation, support a future
network, or provide an independent measurement geometry. But each
answer should be tied to an identifiable technical requirement.

If the only objective is to demonstrate that an atomic clock
can operate near Mars, the proposed mission may be unnecessarily
ambitious. If the objective is to demonstrate an end-to-end
navigation architecture with specific performance targets, the
dedicated spacecraft has a stronger justification.

\textbf{C. The trajectory may not require nuclear electric
	propulsion}

An atomic-clock payload does not intrinsically require an electric
thruster. A clock-equipped spacecraft could reach Mars using a
chemical launch and transfer architecture, potentially with
chemical orbit insertion.

Whether NEP is advantageous depends on spacecraft mass, power,
launch constraints, transfer duration, the desired Mars orbit,
and the propellant mass saved by electric propulsion.

The project must avoid circular reasoning: it should not assume
that an NEP spacecraft is needed merely because the group wants
to design an NEP thruster.

A proper comparison would evaluate the same payload and mission
requirements under at least two propulsion architectures. If
NEP significantly improves payload fraction or enables a
trajectory that the baseline cannot support, the propulsion
case is stronger. If the benefit is small, the mission may still
be useful for navigation research but weak as a justification
for the thruster.

\textbf{D. Orbit insertion and mission duration may dominate
	the design}

A bound Mars orbit requires the spacecraft to arrive with a
suitable velocity and energy. A low-thrust arrival architecture
can be more complicated than simply adding a cruise \(\Delta v\)
to a transfer estimate.

The trajectory must account for the desired arrival orbit, the
available thrust, planetary gravity, and the time at which thrust
is applied. The feasibility of capture or orbit acquisition
depends on the specific trajectory and propulsion design.

For an undergraduate project, this creates a risk of spending
most of the effort on trajectory optimization rather than on
the Hall thruster itself. The group can manage this risk by
selecting a reference orbit, using an established trajectory
tool, and treating detailed orbit design as a bounded subsystem
study.

\textbf{E. The clock payload does not necessarily demand
	large propulsion capability}

A small clock instrument may not create a particularly large
spacecraft mass or power requirement. The required NEP system
could therefore be driven more by transfer duration and spacecraft
architecture than by the payload.

This is not automatically a disadvantage. It does mean, however,
that the team must not use the clock payload as a substitute
for a mass and power budget for the project.

The clock's actual electrical power, thermal control needs,
radio equipment, antenna, data handling, attitude control, and
spacecraft bus must be estimated separately. The thruster should
then be sized around the complete spacecraft, not around the
payload label.

\subparagraph{Assessment}

The Mars-orbit atomic-clock mission is strongest when defined as
a precision navigation or timing demonstrator with a clear
measurement objective. It is less convincing when defined only
as ``put an atomic clock around Mars.''

Its strongest advantages are the direct connection to Mars
operations, an established technical precedent for clock-enabled
deep-space navigation, repeated access to the Martian environment,
and a reasonably clear path from trajectory requirements to
thruster sizing.

Its principal weaknesses are the need for a real navigation
architecture, the possibility that an existing spacecraft could
host the payload, and the fact that NEP must be justified
independently of the clock's scientific value.

% ============================================================
\paragraph{Mission B: Atomic Clock near Mars--Sun L1}
% ============================================================

\subparagraph{Clarifying the location}

In this section, ``L1'' means the first collinear Lagrange point
of the Mars--Sun system, not the Earth--Sun L1 point and not
a point in the Mars--spacecraft system.

The idealized Mars--Sun L1 point lies approximately between Mars
and the Sun. In a rotating three-body model, it is an equilibrium
location in the frame rotating with the two massive bodies.
It is not a location where all forces vanish in an inertial frame.

A spacecraft would generally not be placed at a perfectly fixed
point. A practical mission might use a bounded orbit near L1,
with its actual trajectory and station-keeping needs determined
by the dynamical model and mission requirements.

The precise orbit matters. A mission described only as
``at L1'' is not yet a complete trajectory specification.

\subparagraph{Arguments in favor}

\textbf{A. It provides a different geometric location for
	navigation measurements}

The most defensible argument for the L1 architecture is that
it can occupy a location different from a Mars orbiter. If the
spacecraft is part of a distributed timing and navigation system,
that separation may create measurement geometry that is unavailable
from a Mars-centered orbit alone.

For example, suppose a future network includes multiple Mars
orbiters, surface assets, and spacecraft with known or estimated
states. A clock-equipped spacecraft near Mars--Sun L1 could
potentially provide an additional measurement node. Depending
on the geometry and measurement types, this node might improve
the conditioning of the navigation problem or help distinguish
some state errors from others.

However, spatial separation is not automatically useful. A
measurement contributes most when it adds information that is
not already supplied by other measurements. A second spacecraft
with nearly redundant geometry may add less value than one
positioned to resolve an existing ambiguity.

The claim should therefore be tested through an estimation model.
A suitable comparison would use the same navigation system with
and without the L1 spacecraft, then compare the resulting
state-estimation errors and robustness.

\textbf{B. It creates a more distinctive research question}

The L1 architecture can be intellectually attractive because it
requires the group to consider the interaction between
interplanetary dynamics and a distributed timing system.

The project could examine how spacecraft position, clock stability,
radio-link geometry, and orbit uncertainty influence navigation
performance. It could also compare different bounded orbits near
L1, subject to a clearly defined mission requirement.

This is a stronger research question than merely asking whether
an atomic clock can operate near Mars. It asks whether a specific
spacecraft placement creates a measurable advantage in a larger
system.

For a technically ambitious group, this can lead to a useful
systems-engineering study involving

\begin{itemize}
	\item three-body orbital dynamics;
	\item trajectory design and station-keeping;
	\item clock error and frequency stability;
	\item radio-ranging or Doppler measurement models;
	\item navigation observability and covariance analysis;
	\item sensitivity to spacecraft orbit uncertainty.
\end{itemize}

The main value is the opportunity to formulate and test a
nontrivial systems-level hypothesis. It is not evidence that
the L1 architecture is inherently superior.

\textbf{C. It may support a future distributed infrastructure
	architecture}

A dedicated node near L1 could be considered as part of a future
Mars communications, timing, or navigation infrastructure.
Its usefulness would depend on what the rest of the system
contains.

For instance, the architecture might include multiple spacecraft
with stable clocks, cross-links, radio ranging, known ephemerides,
and a method for distributing time or navigation solutions to
Mars orbiters and surface assets. Under such an architecture,
the L1 node could contribute an additional reference or
measurement path.

This is a plausible concept, but it is a system-level proposition.
It should not be confused with an immediate operational requirement
for Mars exploration.

The group would need to explain why the node belongs near L1
rather than in Mars orbit, in a different solar orbit, or on
another spacecraft already serving the network. The answer
should follow from the geometry and performance analysis.

\subparagraph{Arguments against}

\textbf{A. L1 does not provide an intrinsic atomic-clock
	advantage}

There is no general principle that an atomic clock becomes more
accurate or useful merely because its spacecraft is near L1.

The clock's frequency stability is determined by the clock
technology, environmental conditions, and its implementation.
The mission location affects the system through dynamics,
communications, radiation, thermal conditions, and measurement
geometry, but the location alone does not guarantee improved
timing performance.

The L1 architecture therefore needs an external justification:
a navigation or timing task that benefits from the location.
Without that task, the mission risks being a technically
interesting placement exercise without a sufficiently clear
operational purpose.

\textbf{B. A single clock at L1 does not create a navigation
	network}

The same limitation applies here as in Mars orbit, but it is
particularly important for the L1 concept. A clock-equipped
spacecraft near L1 does not automatically provide position
fixes to Mars orbiters or rovers.

A navigation service requires defined signals and measurements.
Depending on the design, these might include one-way or
two-way radio links, Doppler measurements, ranging, cross-links,
known ephemerides, and clock synchronization or calibration.

If the spacecraft does not communicate with other nodes in
a way that improves their state estimates, its location near
L1 has no operational navigation benefit.

A convincing proposal must answer at least four questions:

\begin{enumerate}
	\item Which spacecraft or surface assets receive a benefit?
	\item Which measurements does the L1 spacecraft provide?
	\item What information does it add beyond existing nodes?
	\item How much does that information improve a stated
	      navigation or timing metric?
\end{enumerate}

If these questions cannot be answered, the L1 mission is not
yet justified as a navigation mission.

\textbf{C. The dynamical environment introduces additional
	trajectory requirements}

A practical L1 mission needs an actual reference trajectory,
not just a nominal equilibrium point. The spacecraft may follow
a bounded orbit around the point, and the trajectory can be
sensitive to injection errors and perturbations.

The exact station-keeping requirements depend on the chosen
orbit, the dynamical model, the mission duration, and acceptable
position error. They cannot be established merely by noting
that L1 is a Lagrange point.

The team must also define how the spacecraft reaches the
operational orbit. A transfer from Earth to the vicinity of
Mars--Sun L1 is not the same problem as a transfer to Mars
orbit. The spacecraft may never need to enter a Mars-centered
bound orbit, but it still needs to reach the target region
with appropriate position and velocity.

For the Hall-thruster project, the trajectory can become a
substantial study in its own right. This may be desirable
if the team wants to explore three-body dynamics and low-thrust
control, but it creates scope risk if the primary deliverable
is a thruster design.

\textbf{D. It may be difficult to demonstrate a direct
	operational need}

Mars orbiters already have clear roles in imaging, communications
relay, atmospheric science, and other Mars operations. An L1
clock spacecraft needs a similarly explicit function.

If the mission's only proposed benefit is improved navigation,
then the group should compare the L1 node with alternatives:
another Mars orbiter, a surface timing node, a different
interplanetary orbit, or an additional clock on an existing
spacecraft.

The L1 architecture becomes compelling only if it produces
a meaningful advantage that these alternatives do not provide
as effectively. Otherwise, it may be more complicated without
being more useful.

\textbf{E. The propulsion benefit is not guaranteed}

The L1 location does not, by itself, imply that Hall propulsion
will be better than conventional propulsion. The answer depends
on the transfer trajectory, spacecraft mass, available power,
flight time, and target-orbit requirements.

A high-specific-impulse system can reduce propellant consumption,
but its lower thrust at fixed power may increase transfer time.
A station-keeping requirement may favor continuous or intermittent
electric propulsion, but that must be demonstrated for the
chosen trajectory.

The project should not assume that an L1 mission is an especially
good NEP application simply because it is unusual. It must
show a quantitative benefit relative to a reasonable baseline.

\subparagraph{Assessment}

The L1 architecture has its strongest case as a research-oriented
distributed-navigation concept. Its potential advantage is
different measurement geometry and a possible role in a broader
infrastructure system.

Its weakness is that these advantages are conditional. Without
a specified network, measurement model, bounded orbit, and
quantitative performance comparison, ``L1'' is a location rather
than a demonstrated mission requirement.

For a project whose main goal is Hall-thruster design, the L1
option should be selected only if the group is prepared to make
the navigation and trajectory analysis part of the central
research question.

% ============================================================
\paragraph{Mars Orbit versus Mars--Sun L1}
% ============================================================

\subparagraph{The fundamental difference}

The Mars-orbit concept begins with a relatively direct operational
context: a spacecraft operates in the Martian environment and
may support navigation, timing, or communications services
associated with Mars.

The L1 concept begins with a different proposition: a spacecraft
is placed at a distinct location in the Mars--Sun dynamical
environment in order to create a potentially useful measurement
or infrastructure node.

The first has a clearer immediate application. The second has
a more conditional systems-level argument.

Neither location inherently produces better navigation. The
relevant quantity is the information available to the navigation
system after accounting for all measurements, errors, and
uncertainties.

\subparagraph{A quantitative way to compare them}

Let the spacecraft state be represented by

\[
	\mathbf{x}
	=
	\begin{bmatrix}
		\mathbf{r} \\
		\mathbf{v} \\
		b          \\
		d
	\end{bmatrix},
\]

where \(\mathbf{r}\) is position, \(\mathbf{v}\) is velocity,
\(b\) is clock bias, and \(d\) is clock drift. A measurement
model can be written abstractly as

\[
	\mathbf{y}
	=
	h(\mathbf{x})+\boldsymbol{\epsilon},
\]

where \(h\) maps the state to predicted measurements and
\(\boldsymbol{\epsilon}\) represents measurement error.

The navigation value of an additional spacecraft depends on
how its measurements constrain the unknown state. In a
linearized estimation framework, the measurement sensitivity
matrix is

\[
	H
	=
	\frac{\partial h}{\partial \mathbf{x}}.
\]

If the measurement-noise covariance is \(R\), a simplified
information contribution is proportional to

\[
	\mathcal{I}
	=
	H^{\mathsf T}R^{-1}H.
\]

This expression is not a complete navigation solution. It
illustrates the key point: a spacecraft is useful when its
measurements add sufficiently independent information relative
to the errors and information already present in the system.

The comparison should therefore hold the clock assumptions,
measurement types, radio-system assumptions, and estimator
constant, then compare the navigation performance obtained
from the two candidate locations.

\subparagraph{Recommended comparison experiment}

The team could formulate the following study:

\begin{enumerate}
	\item Define a baseline navigation network consisting of
	      a ground reference and a specified set of Mars orbiters.

	\item Define a clock stability model and a radio measurement
	      model using published or explicitly assumed parameters.

	\item Simulate the baseline navigation performance without
	      the proposed clock spacecraft.

	\item Add the clock-equipped Mars orbiter and calculate the
	      resulting position and velocity estimation errors.

	\item Repeat the experiment with a spacecraft on a specified
	      bounded orbit near Mars--Sun L1.

	\item Compare accuracy, measurement availability, sensitivity
	      to clock drift, sensitivity to ephemeris error, and robustness
	      to measurement outages.

	\item Perform a separate trajectory and propulsion analysis
	      for each architecture.
\end{enumerate}

The experiment must not assume that the two spacecraft provide
identical measurements from different locations unless that is
physically plausible. The measurement links, visibility, signal
strength, and time dependence should follow from the geometry.

\subparagraph{Decision}

If the team wants the clearest mission justification and the
most direct connection to established Mars operations, Mars
orbit is the stronger baseline.

If the team wants a more novel research problem and is prepared
to investigate navigation geometry, three-body dynamics, and
station-keeping, L1 is a defensible alternative.

The L1 concept should not be selected solely because it sounds
more advanced. It should be selected if a quantitative analysis
shows that the location provides a useful capability that a
Mars-orbit spacecraft does not provide as effectively.

% ============================================================
\paragraph{Mission C: Mars Flyby}
% ============================================================

\subparagraph{Concept and mission objective}

In a Mars-flyby mission, the spacecraft approaches Mars, performs
a close or distant encounter, and then continues along an
outbound trajectory. It does not enter a long-term bound Mars
orbit and does not land.

A flyby could carry a scientific payload, perform a technology
demonstration, collect measurements during the encounter, or
use Mars as a gravity-assist body to alter its heliocentric
trajectory. These are different mission concepts and should
not be treated as interchangeable.

A scientific flyby needs a defined encounter objective. A
gravity-assist mission needs a defined post-encounter trajectory.
A mission that combines both should identify the primary
objective and show how the other supports it.

NASA's descriptions of gravity-assist trajectories explain how
a planetary encounter can change a spacecraft's heliocentric
velocity by exchanging orbital energy and angular momentum
with the planet \cite{NASATrajectories}. A flyby does not
provide unlimited propulsion: the useful change depends on
the encounter geometry and the planet's motion.

\subparagraph{Arguments in favor}

\textbf{A. The mission avoids the requirement to enter
	a long-term Mars orbit}

A flyby does not require the spacecraft to remain bound to
Mars after the encounter. This can simplify the mission
relative to an architecture requiring Mars orbit insertion,
long-term orbit maintenance, or a surface landing.

The spacecraft must still reach the required encounter state.
The arrival position, velocity, closest-approach altitude,
and outbound state are constrained by the desired mission.
However, the spacecraft does not have to satisfy the additional
requirement of remaining in a selected Mars orbit.

For a preliminary study, this provides a relatively clean
mission boundary condition: deliver the spacecraft to a
specified Mars encounter, then determine the resulting trajectory.

This is attractive for an undergraduate group because the
team can focus on the relationship between transfer trajectory,
thruster performance, and total propellant use without designing
a full Mars-orbit operations system.

\textbf{B. A flyby provides a clear trajectory-analysis problem}

A flyby naturally lends itself to a comparison between
propulsion architectures. The team can estimate the transfer
requirements, calculate a feasible encounter geometry, and
compare a chemical baseline against a low-thrust trajectory.

The required analysis may include

\begin{itemize}
	\item launch and departure conditions;
	\item heliocentric transfer trajectory;
	\item thrust magnitude and pointing history;
	\item Mars-relative hyperbolic excess velocity;
	\item closest-approach distance and encounter time;
	\item outbound trajectory and final mission conditions.
\end{itemize}

The flyby can therefore provide a focused way to investigate
whether the selected Hall thruster can deliver the required
integrated velocity change within the available mission time.

\textbf{C. A gravity assist may reduce the propulsion
	burden of a larger mission}

If the spacecraft uses Mars as a gravity-assist body, the
planet can redirect the trajectory and alter the spacecraft's
heliocentric energy. Under suitable geometry, this may reduce
the propulsion needed to reach a later destination or achieve
a desired outbound state.

This can make a flyby useful as part of a larger interplanetary
mission, particularly when the propulsion system has limited
available power or when propellant mass is a major constraint.

The benefit must be assessed against the trajectory that would
be used without the gravity assist. It is not enough to claim
that Mars supplies free \(\Delta v\). The flyby changes the
spacecraft's velocity relative to the Sun through a gravitational
encounter with a moving planet; the result depends on the
incoming asymptote, encounter geometry, and outgoing asymptote.

A flyby can also be combined with electric propulsion. For
example, the thruster may shape the approach and departure
trajectory, while Mars provides an additional gravity-driven
change in heliocentric velocity. Whether this is advantageous
depends on the mission's actual boundary conditions.

\textbf{D. The mission can be used to test a propulsion
	system over a finite operating campaign}

A flyby mission may be designed around a finite cruise and
encounter phase rather than a long period of Mars-local
operations. That can be useful if the project aims to study
thruster operating life, accumulated discharge time, or
propellant consumption during a single transfer.

The thruster requirements remain mission-dependent. A short
mission with a tight arrival-time constraint may favor higher
thrust, while a longer mission with relaxed timing may allow
lower thrust and higher specific impulse. The team should
calculate this trade rather than assuming that a flyby is
automatically easier for electric propulsion.

\textbf{E. The payload can be kept simple}

A flyby can carry a relatively compact instrument or technology
demonstrator. If the main objective is propulsion validation,
the payload may be represented as a specified mass and power
budget.

This allows the team to isolate the propulsion design from
a large surface system, relay network, or long-term orbital
operations architecture. A simplified payload also reduces
the number of spacecraft subsystems that must be estimated.

\subparagraph{Arguments against}

\textbf{A. A flyby may not create a compelling need for NEP}

The absence of a Mars-orbit insertion requirement does not
automatically make a flyby a good electric-propulsion mission.

A conventional chemical trajectory may already accomplish
the flyby with acceptable mass and duration. If the mission
has a relatively modest velocity requirement, a high-power
NEP system may add complexity without a proportionate benefit.

The group must compare the required launch mass, propellant,
power system, trajectory duration, and arrival conditions
against a credible baseline.

A flyby is a useful project when the analysis shows a meaningful
advantage for electric propulsion or when the primary objective
is to demonstrate low-thrust trajectory design. It is weaker
when the mission could be completed more simply by a conventional
architecture and no further benefit is established.

\textbf{B. The encounter is a time- and geometry-constrained
	event}

Unlike an orbiter, a flyby spacecraft has a limited interval
during which it is near Mars. If the science objective depends
on a close encounter, the trajectory must meet the required
closest-approach conditions.

Errors in arrival state can change the encounter geometry,
measurement opportunity, and outbound trajectory. The group
must define acceptable error margins and account for navigation
and maneuver requirements.

A flyby also creates a one-time operational window. There is
no opportunity to repeat the encounter simply by waiting for
the next orbit. A second encounter would require a new trajectory
or a different mission design.

This can make the mission operationally simple after the
encounter but demanding in the approach phase.

\textbf{C. The spacecraft does not provide sustained
	Mars-local infrastructure}

A flyby spacecraft cannot serve as a persistent Mars navigation
node or routine communications relay after it departs. It may
collect useful data during the encounter, but its continued
availability near Mars is limited.

That is a major disadvantage if the intended purpose of the
atomic clock is to support Mars operations over years. The
clock might demonstrate one-way tracking or contribute to
a technology experiment, but it would not be available as
a long-term Mars-local timing reference.

The flyby architecture is therefore poorly matched to a mission
whose main value depends on repeated coverage, sustained
communications, or long-term orbit determination at Mars.

\textbf{D. A gravity assist is not a guaranteed propulsion
	substitute}

The value of a gravity assist depends on the planetary encounter.
A flyby can increase or decrease heliocentric energy, depending
on the geometry. It cannot provide an arbitrary change in
velocity while preserving any desired outgoing trajectory.

The team must solve for the required incoming and outgoing
states. A gravity-assist trajectory that is excellent for one
destination may be unsuitable for another.

Furthermore, a low-thrust spacecraft may spend considerable
time and propellant shaping the encounter trajectory. The
gravity assist must be evaluated as part of the complete
trajectory, not as a standalone source of free performance.

\textbf{E. A flyby can be less informative for long-duration
	thruster operations}

A flyby may be suitable for studying integrated propellant
consumption and finite-duration thruster performance. It is
less directly suited to a project whose objective is to assess
years of operation, repeated starts, long-term erosion, or
the integration of the propulsion system with a persistent
Mars-orbit service.

These effects can still be studied using a flyby mission,
but the connection is less direct. If the group wants to
evaluate long operational lifetime, a long-duration orbital
or infrastructure mission may provide a clearer requirements
basis.

\subparagraph{Assessment}

A Mars flyby is the cleanest of these options for a focused
trajectory and encounter study. It avoids long-term Mars-orbit
operations and surface landing, and it gives the team a clear
set of encounter conditions to target.

Its primary weakness is that it offers little sustained Mars-local
utility. It is also not automatically a strong NEP application:
the mission must demonstrate that electric propulsion provides
a meaningful benefit over a conventional baseline.

The flyby is most appropriate if the group prioritizes a bounded
propulsion and trajectory problem rather than a persistent
Mars navigation or communications capability.

% ============================================================
\paragraph{Mission D: Mars Rover}
% ============================================================

\subparagraph{Concept and mission objective}

A Mars rover mission delivers a mobile surface platform to
Mars to conduct geological, environmental, astrobiological,
or other surface investigations. The rover may carry scientific
instruments, communications equipment, autonomous navigation
systems, and power systems needed for surface operation.

The rover is an obvious mission concept because its scientific
purpose is easy to explain. A rover can travel to a selected
region, investigate local terrain, and return data through
Mars orbiters or a direct-to-Earth communications link.

However, a rover mission is a coupled transportation and
surface-system problem. It requires not only interplanetary
transport but also a viable Mars arrival and landing architecture.

\subparagraph{Arguments in favor}

\textbf{A. The scientific and operational objective is concrete}

A rover has a direct, understandable purpose: reach a specified
location and perform surface science or exploration.

That gives the project a useful basis for a payload mass budget,
mission duration, surface power estimate, and communications
requirement. The team can specify a particular science region
or surface task rather than just a general spacecraft destination.

A concrete payload also makes it easier to define the spacecraft
configuration. The rover, landing system, cruise stage, and
communications equipment each contribute mass and operational
requirements.

For mission-level propulsion sizing, this can produce a meaningful
payload-driven problem: the propulsion system must transport
the complete vehicle, including the hardware needed to land
and operate on Mars.

\textbf{B. Electric propulsion may reduce the mass of
	the interplanetary transportation stage}

For a sufficiently long transfer, high-specific-impulse
electric propulsion may reduce the propellant required for
interplanetary transportation relative to a chemical system
with comparable mission constraints.

If propellant mass is reduced, some of the saved mass may
be available for the rover, scientific instruments, landing
hardware, or power system. The actual result depends on the
spacecraft architecture, the available electrical power, the
transfer duration, and the required arrival conditions.

This is an important distinction: electric propulsion does
not automatically increase payload mass. The power system
has its own mass, and lower thrust can lengthen the transfer.
The correct comparison must include the entire propulsion
and power subsystem rather than just the propellant tank.

A useful design question is whether a nuclear-electric
transport stage can deliver a specified rover and landing
system to Mars with a better payload fraction, lower total
mass, or improved mission flexibility than the selected
baseline architecture.

\textbf{C. The rover creates a real mass and power budget}

Unlike a purely abstract payload, a rover has identifiable
subsystems. The mission analysis can assign masses to the
rover chassis, wheels, instruments, onboard computing,
communications equipment, thermal control, and surface power.

The landing system adds a separate mass contribution. Depending
on the architecture, it may include an aeroshell, thermal
protection, parachute, propulsion for terminal descent, landing
legs, or other hardware.

These contributions make the total spacecraft mass more
meaningful and can constrain the propulsion design.

The group should distinguish between

\[
	m_{\mathrm{spacecraft}}
	=
	m_{\mathrm{cruise}}
	+
	m_{\mathrm{rover}}
	+
	m_{\mathrm{landing}}
	+
	m_{\mathrm{payload}}
	+
	m_{\mathrm{propellant}},
\]

where the exact categories must be defined to avoid double-counting.
The equation is a bookkeeping framework, not a substitute for
a detailed mass budget.

\textbf{D. The mission offers a clear case for relay communications}

Mars surface assets often rely on orbiters to relay substantial
amounts of data to Earth. NASA describes the Mars Relay Network
as an important part of the communications architecture for
surface missions \cite{MarsRelayNetwork}.

A rover mission can therefore be paired with a broader
communications or navigation concept. For example, the rover
might benefit from an orbital network that provides command
and telemetry relay, and a future navigation system could
potentially improve knowledge of spacecraft or rover states.

An atomic clock could be studied as a separate component of
that system, but it should not be assumed that the rover itself
requires a dedicated clock-equipped spacecraft. The group
would need to identify the specific measurement or operational
benefit.

\textbf{E. It has substantial engineering depth}

A rover mission provides opportunities to study interactions
between propulsion, spacecraft power, payload mass, thermal
control, and landing architecture.

For an undergraduate team, this can be an advantage if the
goal is a multidisciplinary spacecraft design project. The
project could include a top-level mass budget, an interplanetary
trajectory, a propulsion trade, and a preliminary system
architecture.

It also creates a clear division of responsibilities. One
subgroup can develop the Hall thruster and its operating map,
another can estimate the transfer trajectory, and another
can establish the rover and landing-system mass budget.

This division is only useful if the interfaces are defined
early. Otherwise, each subsystem may evolve around incompatible
assumptions.

\subparagraph{Arguments against}

\textbf{A. Atmospheric entry, descent, and landing are
	separate major engineering problems}

The most important drawback is that a rover must survive
arrival at Mars and reach the surface safely. This introduces
a mission phase that is not solved by the interplanetary
propulsion system.

The landing system must manage entry heating, aerodynamic
loads, descent dynamics, and terminal landing conditions.
The detailed architecture depends on rover mass, landing
site elevation, atmospheric assumptions, terrain, and
acceptable risk.

An electric thruster used during interplanetary cruise does
not automatically solve these problems. The team should not
conflate a low-thrust transfer trajectory with the atmospheric
entry and landing design.

If the main project deliverable is a Hall thruster, designing
a rover and its full landing system could consume substantial
effort without improving the thruster model. The landing
system should therefore be treated as a defined payload
mass and interface unless entry and landing are explicitly
part of the project.

\textbf{B. The rover's surface power system is not the
	same as the spacecraft's electric-propulsion power system}

A rover needs power after landing. A Hall thruster requires
substantial electrical power during operation, along with
power-processing equipment and associated thermal management.

These requirements occur in different mission phases. The
interplanetary spacecraft may use a nuclear-electric power
system, while the rover may use a separate surface power
architecture. The team must decide whether the nuclear power
system is used only for transport, remains with the lander,
or powers a larger integrated architecture.

Each choice changes the mass budget and mission design.
For example, carrying a large power system through landing
may introduce additional mass and complexity, while separating
the cruise stage from the landed system changes the spacecraft
interfaces and operational concept.

The group should not simply add a rover payload to a Hall
thruster and assume that the complete vehicle is a consistent
design.

\textbf{C. The mission may become a landing-system project
	rather than a thruster project}

A rover is a compelling exploration platform, but its
requirements may be dominated by surface mobility, scientific
instruments, thermal survival, landing risk, and communications.

If the team tries to design all these elements in detail,
the Hall-thruster study may become a relatively small part
of the overall work. This is a scope-management problem,
not a scientific argument against rover missions.

A reasonable boundary is to define the rover and landing
system through a reference mass, power requirement, and
interface specification. The propulsion team can then
evaluate the transfer stage required to deliver that payload.

\textbf{D. The rover does not inherently require an
	atomic clock}

A rover mission may benefit from accurate timing and navigation,
but the presence of a rover does not establish the need for
a dedicated atomic-clock spacecraft.

The navigation architecture could rely on orbiters, ground
tracking, onboard sensors, local mapping, and other techniques.
Whether a precision clock adds useful capability depends on
the task and the measurement architecture.

If the mission includes an atomic clock, the proposal should
explain whether it supports interplanetary navigation, Mars-orbit
navigation, rover localization, or a broader network. These
are distinct functions and require different performance
analyses.

A clock-equipped Mars orbiter might be a useful addition
to a rover-support architecture, but it should be justified
as a separate subsystem rather than treated as an automatic
consequence of selecting a rover.

\textbf{E. The mass benefit of NEP must be evaluated across
	the complete vehicle}

A rover mission can make high specific impulse attractive
because the payload and landing system may be substantial.
However, the power system and electric-propulsion hardware
also contribute mass.

A fair comparison should include at least

\begin{itemize}
	\item thruster and power-processing-unit mass;
	\item nuclear-electric power-system mass;
	\item propellant and tank mass;
	\item structural and thermal mass;
	\item cruise-stage mass;
	\item rover and scientific payload mass;
	\item landing-system mass;
	\item any additional hardware needed for the selected
	      transfer and arrival architecture.
\end{itemize}

The result may favor NEP, chemical propulsion, or a hybrid
architecture. The answer depends on the mission and cannot
be inferred from specific impulse alone.

\subsection{Assessment}

A rover is the strongest candidate if the group's central
objective is an exploration mission with a concrete surface
payload. It offers a clear science case and a meaningful
spacecraft mass budget.

It is less suitable if the goal is to isolate Hall-thruster
design and computational validation. The landing system and
surface operations create a large set of requirements that
must either be handled explicitly or bounded as external
interfaces.

For a propulsion-focused project, the rover should be treated
as a payload-delivery mission. The team should avoid expanding
the scope to include full rover, entry, descent, landing,
and surface-power design unless those are explicit project
deliverables.

% ============================================================
\paragraph{Cross-Mission Comparison}
% ============================================================

\subparagraph{Qualitative comparison}

The following table is a qualitative engineering assessment.
It is not a numerical ranking derived from a trajectory simulation.
The terms ``high'' and ``low'' refer to the relative clarity
or complexity of the mission concept for the stated project.

\begin{longtable}{
		p{0.17\textwidth}
		p{0.19\textwidth}
		p{0.19\textwidth}
		p{0.19\textwidth}
		p{0.19\textwidth}
	}
	\toprule
	\textbf{Criterion}
	 & \textbf{Mars-orbit clock}
	 & \textbf{L1 clock}
	 & \textbf{Mars flyby}
	 & \textbf{Mars rover}
	\\
	\midrule
	\endfirsthead

	\toprule
	\textbf{Criterion}
	 & \textbf{Mars-orbit clock}
	 & \textbf{L1 clock}
	 & \textbf{Mars flyby}
	 & \textbf{Mars rover}
	\\
	\midrule
	\endhead

	Immediate mission purpose
	 & Strong if tied to navigation or timing
	 & Conditional on a defined network role
	 & Strong if a specific encounter or science objective is selected
	 & Strong and easy to communicate
	\\

	Connection to existing Mars operations
	 & Direct
	 & Indirect unless integrated into a network
	 & Limited after the encounter
	 & Direct through surface exploration and relay
	\\

	Trajectory complexity
	 & Moderate to high; depends on arrival orbit
	 & High; target-region dynamics and bounded orbit must be specified
	 & Moderate to high; encounter conditions are critical
	 & High when the complete arrival and landing architecture is included
	\\

	Need for long-term Mars operations
	 & Yes, if the clock is intended as a persistent Mars node
	 & No Mars orbit is required, but sustained operations may still be intended
	 & No
	 & Yes, for surface exploration
	\\

	Navigation-system research potential
	 & Strong and directly motivated
	 & Potentially strong but conditional
	 & Limited for persistent Mars navigation
	 & Depends on whether navigation is a central payload objective
	\\

	Landing-system complexity
	 & None
	 & None
	 & None
	 & Major additional subsystem
	\\

	Scope risk for a thruster project
	 & Moderate
	 & High if navigation and three-body dynamics become dominant
	 & Low to moderate for a bounded flyby study
	 & High if landing and surface systems are included
	\\

	Clarity of NEP justification
	 & Must be shown with trajectory and mass comparison
	 & Must be shown with trajectory and mass comparison
	 & Must be shown; chemical baseline may be competitive
	 & Potentially strong for a large payload, but system mass must be counted
	\\

	Best project emphasis
	 & Navigation and propulsion co-design
	 & Distributed navigation and three-body dynamics
	 & Low-thrust trajectory and encounter design
	 & Payload delivery and multidisciplinary systems engineering
	\\
	\bottomrule
\end{longtable}

\subparagraph{The four missions answer different questions}

The missions should not be ranked as if they were interchangeable
payloads on the same trajectory.

The Mars-orbit clock asks whether an onboard clock and a
Mars-based spacecraft can improve navigation or timing in
the Martian environment.

The L1 clock asks whether a spacecraft near a different
dynamical location can contribute useful information to a
distributed navigation or timing system.

The flyby asks whether the propulsion system can meet a
finite-duration interplanetary encounter requirement, possibly
including a useful gravity assist.

The rover asks whether the propulsion and transportation
architecture can deliver a comparatively complex surface
payload while meeting the required arrival and landing
conditions.

The preferred mission depends on which of these questions
the group actually wants to answer.

% ============================================================
\paragraph{Recommended Quantitative Evaluation}
% ============================================================

\subparagraph{Step 1: Define the mission requirements}

Before selecting a thruster operating point, establish the
requirements that all four architectures must satisfy or
explicitly identify which requirements differ.

At minimum, specify

\begin{itemize}
	\item launch date or launch-window assumptions;
	\item initial spacecraft mass and payload mass;
	\item destination or encounter conditions;
	\item allowable transfer duration;
	\item required operational lifetime;
	\item available electrical power;
	\item arrival or encounter velocity constraints;
	\item communication and navigation objectives;
	\item propellant and thermal constraints.
\end{itemize}

A mission concept without these parameters cannot support a
defensible thruster selection. If exact values are not yet
available, use a documented range and conduct sensitivity
analysis.

\subparagraph{Step 2: Build a baseline mass budget}

For each mission, separate the payload, spacecraft bus,
power system, propulsion system, and propellant. Avoid
double-counting the cruise stage or landing system.

A useful initial accounting identity is

\[
	m_0 =
	m_{\mathrm{payload}}
	+
	m_{\mathrm{bus}}
	+
	m_{\mathrm{power}}
	+
	m_{\mathrm{thruster}}
	+
	m_{\mathrm{PPU}}
	+
	m_{\mathrm{tank}}
	+
	m_p
	+
	m_{\mathrm{other}},
\]

where \(m_{\mathrm{other}}\) represents any additional
subsystems not listed separately. The definition of each
term must be consistent across the comparison.

For the rover case, the rover and landing system should be
represented separately. For the clock missions, the clock,
radio hardware, and any communications or navigation equipment
should be included in the payload or bus budget.

The team should avoid assigning a highly precise value to
a subsystem before the architecture is defined. Early mass
budgets should use credible estimates and uncertainty ranges.

\subparagraph{Step 3: Compute first-order propulsion requirements}

For a chosen specific impulse and efficiency, estimate the
thrust available at the selected electrical power:

\[
	T \simeq
	\frac{2\eta_t P_{\mathrm{in}}}
	{g_0 I_{\mathrm{sp}}}.
\]

The corresponding mass-flow rate is

\[
	\dot{m} =
	\frac{T}{g_0 I_{\mathrm{sp}}}.
\]

The approximate propellant consumed during a period of
continuous operation is

\[
	m_p \simeq \dot{m}t_{\mathrm{on}}.
\]

This last expression is only a first-order estimate when
the operating point is fixed. A trajectory with variable
thrust, variable power, or changing specific impulse requires
integration over time:

\[
	m_p =
	\int_{t_0}^{t_f} \dot{m}(t)\,dt.
\]

The group should calculate several operating points rather
than select a single arbitrary value. A useful trade study
would vary specific impulse, power, and efficiency, then
determine which combinations are compatible with the
trajectory and lifetime requirements.

\subparagraph{Step 4: Model the trajectory at an appropriate fidelity}

A first-order impulsive transfer estimate can help establish
whether a mission is plausible, but a Hall-thruster mission
ultimately requires low-thrust trajectory analysis.

The trajectory model should include the relevant gravitational
environment, spacecraft mass variation, thrust magnitude,
thrust direction, and mission boundary conditions.

For Mars orbit, the arrival condition should specify the
target orbit or a defined orbit-insertion objective.

For L1, the target should be a specified bounded orbit or
other explicitly defined operational trajectory near the
Mars--Sun L1 region.

For a flyby, the model should specify encounter time,
closest-approach distance, incoming and outgoing states,
and any post-encounter objective.

For a rover, the transfer should deliver the defined lander
and rover mass to the required Mars arrival state. The entry,
descent, and landing system can remain outside the detailed
thruster model if its mass and interface requirements are
explicitly represented.

\subparagraph{Step 5: Compare against a credible alternative}

Each mission should be evaluated against a baseline that
does not assume the proposed thruster is automatically
the best option.

Possible comparisons include

\begin{itemize}
	\item chemical transfer versus electric transfer;
	\item chemical orbit insertion versus electric orbit acquisition;
	\item electric propulsion with different power levels;
	\item different specific-impulse operating points;
	\item different mission durations;
	\item different payload masses;
	\item for the clock concepts, Mars orbit versus L1 using
	      the same navigation-performance assumptions.
\end{itemize}

The team should compare complete system outcomes, including
power-system mass and transfer time, rather than compare
only propellant mass or specific impulse.

\subparagraph{Step 6: Tie the mission to the thruster design}

Once the mission operating envelope is defined, it can
inform the thruster's preliminary design.

Relevant design variables include

\begin{itemize}
	\item discharge voltage and current;
	\item channel inner and outer radii;
	\item channel length;
	\item magnetic-field topology and field strength;
	\item propellant injection and distribution;
	\item cathode placement and coupling;
	\item thermal loads and material limits;
	\item expected operating duration and erosion exposure.
\end{itemize}

These variables should not be selected independently of the
required operating point. For example, the mission-level
power and thrust target may constrain the discharge power,
while the chosen propellant and desired specific impulse
influence the voltage and mass-flow regime.

A simulation of the entire thruster should be treated as a
separate validation hierarchy. A spacecraft trajectory model
does not validate the plasma physics, and a plasma simulation
does not establish mission feasibility by itself.

A defensible workflow is

\[
	\text{mission requirements}
	\rightarrow
	\text{trajectory}
	\rightarrow
	\text{thruster operating point}
	\rightarrow
	\text{thruster geometry and physics}
	\rightarrow
	\text{performance validation}.
\]

The chain must also be iterated in the reverse direction:
if the thruster cannot deliver the assumed performance,
the trajectory and mission mass budget must be updated.

% ============================================================
\paragraph{Final Recommendation}
% ============================================================

\subparagraph{If the priority is a clear mission justification}

Choose the \textbf{atomic clock in Mars orbit} as the baseline
concept, provided that the mission is defined as a navigation
or timing demonstrator rather than merely as a clock placed
near Mars.

It has a direct operational context, an identifiable relationship
to deep-space navigation research, and a straightforward
way to define a bounded mission architecture. The team can
compare conventional and electric-propulsion transfer options,
then assess whether the resulting spacecraft can meet the
required orbit and operational lifetime.

The principal unresolved question is whether a dedicated
NEP spacecraft offers enough advantage over hosting the
clock on another platform. That question should be addressed
through the mass, power, and trajectory comparison.

\subparagraph{If the priority is a more novel navigation study}

Choose the \textbf{Mars--Sun L1 clock} only if the group is
willing to investigate the distributed navigation system
in detail.

The proposal should identify the other nodes, the measurements
made by the L1 spacecraft, the target bounded orbit, and the
metric by which the new node improves system performance.
The analysis should compare the L1 architecture against
a Mars-orbit alternative using consistent clock and radio
assumptions.

Without that analysis, L1 is a less defensible choice because
the location itself does not establish a useful navigation
service.

\subparagraph{If the priority is a focused trajectory project}

Choose the \textbf{Mars flyby} if the primary deliverable is
a well-bounded low-thrust trajectory and encounter study.

It avoids landing-system design and long-term Mars-orbit
operations. It can support a clear comparison between
propulsion architectures, but the team must show that the
electric-propulsion option provides a meaningful benefit
for the chosen encounter and post-encounter objective.

\subparagraph{If the priority is an exploration mission}

Choose the \textbf{Mars rover} if the group wants a
multidisciplinary payload-delivery mission with an obvious
scientific purpose.

For a thruster-focused project, however, define the rover
and landing system as an externally specified payload and
interface. Do not allow full entry, descent, landing, and
surface-system design to dominate the propulsion work unless
those subjects are part of the intended scope.

\subparagraph{Overall conclusion}

For a group whose main deliverable is the design and
computational evaluation of a Hall thruster, the recommended
starting point is a Mars-orbit atomic-clock mission with a
clearly stated navigation objective. It offers a useful
balance between mission relevance, technical specificity,
and manageable scope.

The L1 clock is a stronger choice only if a distributed
navigation analysis can demonstrate a measurable benefit
from its location. The flyby is the cleanest alternative
for a propulsion-and-trajectory-focused study. The rover
is the strongest surface-exploration concept, but it carries
the greatest risk of expanding the project beyond thruster
design.

The final choice should be made only after a preliminary
mass budget and trajectory comparison. In particular, the
team should establish whether NEP improves payload fraction,
transfer feasibility, or mission flexibility relative to
a credible alternative. That result, rather than the
apparent sophistication of the destination, should determine
which mission becomes the baseline.

% ============================================================
\begin{thebibliography}{9}

	\bibitem{DSACNavigation}
	A. S. Ely et al.,
	\textit{Advancing Navigation, Timing, and Science with the
		Deep Space Atomic Clock},
	NASA Technical Reports Server, 2016.
	\url{https://ntrs.nasa.gov/citations/20160008206}

	\bibitem{DSACParadigm}
	\textit{The Deep Space Atomic Clock: Ushering in a New Paradigm
		for Radio Navigation and Science},
	NASA Technical Reports Server.
	\url{https://ntrs.nasa.gov/citations/20150007332}

	\bibitem{MarsRelayNetwork}
	NASA Science,
	\textit{Mars Relay Network}.
	\url{https://science.nasa.gov/mars/mars-relay-network/}

	\bibitem{NASATrajectories}
	NASA Science,
	\textit{Basics of Space Flight: Trajectories}.
	\url{https://science.nasa.gov/learn/basics-of-space-flight/chapter4-1/}

	\bibitem{MarsRelayTechnical}
	NASA Technical Reports Server,
	\textit{A Martian Telecommunications Network: UHF Relay Support
		of the Mars Exploration Rovers by the Mars Global Surveyor,
		Mars Odyssey, and Mars Express Orbiters}.
	\url{https://ntrs.nasa.gov/search.jsp?R=20070023754}

	\bibitem{DSACOnboard}
	NASA Technical Reports Server,
	\textit{Preliminary Investigation of Onboard Orbit Determination
		Using Deep Space Atomic Clock Based Radio Tracking}.
	\url{https://ntrs.nasa.gov/citations/20190025513}

\end{thebibliography}



\textbf{Outline of plan}
\begin{enumerate}
	\item Find parameters
	\item Baseline operating point(discharge voltage, flow rate, etc)
	\item Find secondary performance via calculations
	\item Design vague physical architecture
	\item Optimize for size of channel and magnetic circut
	\item Design propellant delivery
	\item Simulate and further optimize
	\item Integrate models
	\item Write report
\end{enumerate}

Once we have the basic goal and such
\begin{itemize}
  \item Plasma physics and discharge
  \item Electromagnetic coil
  \item Propellant system
  \item Simulate and optimize
\end{itemize}

\section{Design}


\section{Optimization}



\section{Teachings}

\subsection{Numerics}

\textbf{Preliminary Workshop}

I didn't use this for planning but I will probably write notes on how it went and next steps.

\textbf{CFL/finite Difference}

Didn't use this for planning but may not evaluation and next steps

\textbf{*MHD}


\subsection{Plasma Physics}


\textbf{First Go}

\begin{itemize}
	\item What is plasma
	\item Lorentz Force
	\item Parameters
	\item Types of representation's
	\item MHD in more detail
	\item How plasma Produces Thrust
	\item Plasma in relation to our Thruster
\end{itemize}
